\subsection*{Primary geometry across the conductor boundary}
The principal boundary theorem is a classification of an entire
Hilbert-function class, rather than a further isolated local example.
Every complex local algebra with Hilbert function $(1,2,2)$ is
isomorphic to $\C[x,y]/((x,y)^3,y^2)$ or
$\C[x,y]/((x,y)^3,xy)$. For all its unital three-planes, in every
degree at least two, Theorem~\ref{thm:loewy-primary} determines the
full zeroth-Fitting scheme and every associated stratum, including
the extreme-corank boundary. On a smooth frame cover put
$\delta=ad-bc$, $P=(a,c)$ and $Q=(b,d)$. In the square case the ideal
is $\delta^2P^2=(\delta^2)\cap P^4$. In the two-factor case it is
$\delta^2(ab,ad+bc,cd)$, with embedded primes $P,Q,P+Q$ and the
explicit primary decomposition \eqref{eq:loewy-two-factor-primary}.
Both schemes have a normal reduced determinant divisor, generic
multiplicity two, and nilradical index four.

The factor lines of the quadratic relation index the embedded
surfaces. A double factor gives a surface contained entirely in the
subalgebra locus; distinct factors give two surfaces meeting that
locus and an additional embedded point at the quadratic layer.
Theorem~\ref{thm:loewy-conductor-boundary} computes the conductor
kernel on these strata, over their actual coefficient rings. This
is the boundary information not supplied by the corank-one incidence
isomorphism alone. Proposition~\ref{prop:loewy-collision} places
both types in one finite flat algebra family, with an exact formula
for the full Fitting ideal before specialization.

Three levels of statement should be distinguished. For an arbitrary
finite algebra, Theorem~\ref{thm:full-codimension-two} is a global
presentation and an incidence theorem. For cube-zero augmentation
ideals, Theorem~\ref{thm:loewy-factorization} factors the full
multiplication ideal through the quadratic multiplication tensor.
For Hilbert function $(1,2,2)$, Theorem~\ref{thm:loewy-primary}
converts this factorization into a complete primary classification.
The connected higher-contact family of
Theorem~\ref{thm:embedded-multigenerator} supplies a complementary
phenomenon: unbounded nilpotent complexity. Corollary~\ref{cor:loewy-global}
transports the complete length-five classification to section
multiplication on the projective plane.

Finally, Proposition~\ref{prop:codimension-sharpness} proves optimality
of the stabilization bound in every codimension and every generating
rank at least two. Its elementary field-level interpretation is
separated from the coefficient-ring assertion. The arguments in
Sections~\ref{sec:quadratic-fibre-functor} and
\ref{sec:relative-cohomology-details} give the full functorial and
cohomological details needed for arbitrary, possibly nonreduced,
base change.
